% GENERATED FILE. Edit canonical lessons, then run npm run generate:books.
% canonical-source-hash: fnv1a64:698f16d0730b8c3b
% canonical-lessons: RU-C50-wall, RU-C50-roof, RU-C50-garden, RU-C50-bathroom, RU-C50-clock

\chapter{Wall, Roof, Garden, Bathroom, Clock}
\label{ch:ru-wall-roof-garden-bathroom-clock}
\hlchaptermodality{50}

\begin{chapteropening}
\textbf{By the end of this chapter:} \emph{I can name a wall, a roof, a garden, the bathroom and a clock.}

Complete the last lesson of chapter 50: I can name a wall, a roof, a garden, the bathroom and a clock.
\end{chapteropening}

\section[\texorpdfstring{\ru{стена} (stená)}{стена (stená)}]{\ru{стена} (stená) — a wall}

\label{lesson:RU-C50-wall}



\begin{quote}
Before the new one: say the Russian for a lamp, then the Russian for a key.
\end{quote}

\subsection*{You'll want to know: \ru{стена}}
\textbf{\ru{стена}} (\emph{stená}) — \textquotedblleft{}a wall\textquotedblright{}.

\begin{cousinweb}
\textbf{\ru{стена}} $\leftarrow$ Proto-Slavic \emph{*stěna} — inherited. Vasmer relates it to Gothic \emph{stains}, \textquotedblleft{}stone\textquotedblright{}, so English \emph{stone} is a cousin, and he holds that a loan from Germanic cannot be proved.
\end{cousinweb}

\begin{grammarlens}[title={What you've built}]
One more word for this chapter.
\end{grammarlens}

\subsection*{Your turn}
\begin{itemize}
\raggedright
  \item \emph{Say it:} \emph{stená}
  \item \emph{Say it:} \emph{stená}, once more
  \item \emph{Say it:} say \emph{stená}
  \item \emph{Recall:} say the Russian for a lamp, then the Russian for a key, then say \emph{stená} again
\end{itemize}

\begin{tcolorbox}[breakable,colback=teal!4,colframe=teal!35!black,title={Before you move on}]
What does \ru{стена} mean? (\textquotedblleft{}A wall\textquotedblright{}.) Say it once more.
\end{tcolorbox}
\section[\texorpdfstring{\ru{крыша} (krýsha)}{крыша (krýsha)}]{\ru{крыша} (krýsha) — a roof}

\label{lesson:RU-C50-roof}



\begin{quote}
Before the new one: say the Russian for a key, then the Russian for a wall.
\end{quote}

\subsection*{You'll want to know: \ru{крыша}}
\textbf{\ru{крыша}} (\emph{krýsha}) — \textquotedblleft{}a roof\textquotedblright{}.

\begin{cousinweb}
\textbf{\ru{крыша}} is built on \emph{\ru{крыть}}, \textquotedblleft{}to cover\textquotedblright{}, with an \emph{-s-} that became \emph{\ru{ш}} (Vasmer): the roof is the cover. The same verb is in \emph{\ru{открывать}}, \textquotedblleft{}to open\textquotedblright{}, and \emph{\ru{закрывать}}, \textquotedblleft{}to close\textquotedblright{}.
\end{cousinweb}

\begin{grammarlens}[title={What you've built}]
One more word for this chapter.
\end{grammarlens}

\subsection*{Your turn}
\begin{itemize}
\raggedright
  \item \emph{Say it:} \emph{krýsha}
  \item \emph{Say it:} \emph{krýsha}, once more
  \item \emph{Say it:} say \emph{krýsha}
  \item \emph{Recall:} say the Russian for a key, then the Russian for a wall, then say \emph{krýsha} again
\end{itemize}

\begin{tcolorbox}[breakable,colback=teal!4,colframe=teal!35!black,title={Before you move on}]
What does \ru{крыша} mean? (\textquotedblleft{}A roof\textquotedblright{}.) Say it once more.
\end{tcolorbox}
\section[\texorpdfstring{\ru{сад} (sad)}{сад (sad)}]{\ru{сад} (sad) — a garden}

\label{lesson:RU-C50-garden}



\begin{quote}
Before the new one: say the Russian for a wall, then the Russian for a roof.
\end{quote}

\subsection*{You'll want to know: \ru{сад}}
\textbf{\ru{сад}} (\emph{sad}) — \textquotedblleft{}a garden\textquotedblright{}.

\begin{cousinweb}
\textbf{\ru{сад}} $\leftarrow$ Proto-Slavic \emph{*sadŭ}, \textquotedblleft{}a planting\textquotedblright{} — inherited, from \emph{\ru{садить}}, \textquotedblleft{}to plant\textquotedblright{}, literally \textquotedblleft{}to make sit\textquotedblright{}, the partner of \emph{\ru{сидеть}}, \textquotedblleft{}to sit\textquotedblright{} (Vasmer). Old Russian \emph{sadŭ} also meant \textquotedblleft{}a tree, a grove\textquotedblright{}.
\end{cousinweb}

\begin{grammarlens}[title={What you've built}]
One more word for this chapter.
\end{grammarlens}

\subsection*{Your turn}
\begin{itemize}
\raggedright
  \item \emph{Say it:} \emph{sad}
  \item \emph{Say it:} \emph{sad}, once more
  \item \emph{Say it:} say \emph{sad}
  \item \emph{Recall:} say the Russian for a wall, then the Russian for a roof, then say \emph{sad} again
\end{itemize}

\begin{tcolorbox}[breakable,colback=teal!4,colframe=teal!35!black,title={Before you move on}]
What does \ru{сад} mean? (\textquotedblleft{}A garden\textquotedblright{}.) Say it once more.
\end{tcolorbox}
\section[\texorpdfstring{\ru{ванная} (vánnaya)}{ванная (vánnaya)}]{\ru{ванная} (vánnaya) — the bathroom}

\label{lesson:RU-C50-bathroom}



\begin{quote}
Before the new one: say the Russian for a roof, then the Russian for a garden.
\end{quote}

\subsection*{You'll want to know: \ru{ванная}}
\textbf{\ru{ванная}} (\emph{vánnaya}) — \textquotedblleft{}the bathroom\textquotedblright{}.

\begin{cousinweb}
\textbf{\ru{ванная}} is short for \emph{\ru{ванная} \ru{комната}}, \textquotedblleft{}the bath room\textquotedblright{}: an adjective from \emph{\ru{ванна}}, \textquotedblleft{}a bathtub\textquotedblright{}, standing alone. \emph{\ru{ванна}} came from German \emph{Wanne} or French \emph{vanne} (Vasmer); Trubachev adds that it is recorded by 1697.
\end{cousinweb}

\begin{grammarlens}[title={What you've built}]
One more word for this chapter.
\end{grammarlens}

\subsection*{Your turn}
\begin{itemize}
\raggedright
  \item \emph{Say it:} \emph{vánnaya}
  \item \emph{Say it:} \emph{vánnaya}, once more
  \item \emph{Say it:} say \emph{vánnaya}
  \item \emph{Recall:} say the Russian for a roof, then the Russian for a garden, then say \emph{vánnaya} again
\end{itemize}

\begin{tcolorbox}[breakable,colback=teal!4,colframe=teal!35!black,title={Before you move on}]
What does \ru{ванная} mean? (\textquotedblleft{}The bathroom\textquotedblright{}.) Say it once more.
\end{tcolorbox}
\section[\texorpdfstring{\ru{часы} (chasý)}{часы (chasý)}]{\ru{часы} (chasý) — a clock, a watch}

\label{lesson:RU-C50-clock}



\begin{quote}
Before the new one: say the Russian for a garden, then the Russian for the bathroom.
\end{quote}

\subsection*{You'll want to know: \ru{часы}}
\textbf{\ru{часы}} (\emph{chasý}) — \textquotedblleft{}a clock, a watch\textquotedblright{}. Always plural.

\begin{cousinweb}
\textbf{\ru{часы}} is the plural of \emph{\ru{час}}, \textquotedblleft{}an hour\textquotedblright{}: a clock is \textquotedblleft{}the hours\textquotedblright{}. \emph{\ru{час}} $\leftarrow$ Proto-Slavic \emph{*časŭ}, \textquotedblleft{}time, a while\textquotedblright{}, a cousin of Old Prussian \emph{kīsman}, \textquotedblleft{}time\textquotedblright{}, and Albanian \emph{kohë}, \textquotedblleft{}time, weather\textquotedblright{} (Vasmer).
\end{cousinweb}

\begin{grammarlens}[title={What you've built}]
One more word for this chapter.
\end{grammarlens}

\subsection*{Your turn}
\begin{itemize}
\raggedright
  \item \emph{Say it:} \emph{chasý}
  \item \emph{Say it:} \emph{chasý}, once more
  \item \emph{Say it:} say \emph{chasý}
  \item \emph{Recall:} say the Russian for a garden, then the Russian for the bathroom, then say \emph{chasý} again
\end{itemize}

\begin{tcolorbox}[breakable,colback=teal!4,colframe=teal!35!black,title={Before you move on}]
What does \ru{часы} mean? (\textquotedblleft{}A clock, a watch\textquotedblright{}.) Say it once more.
\end{tcolorbox}
